\subsection{cycle}\label{cycle}
The \verb|cycle| module represents the temporal backbone of the slot-based
radio protocol. A complete \emph{cycle} is divided into \verb|SLOT_CNT| slots,
and each slot in turn into \verb|SUB_SLOT_CNT| sub-slots. With the default
values from \verb|system_definitions.h| (Listing \ref{cycle_defines}) this
gives $16$ slots of $8$ sub-slots each, i.e.\ $128$ sub-slots per cycle. Each
device owns its own odd slot ($1, 3, 5, \dots, \mathtt{f}$) and may transmit
only in it; the remaining slots are used to receive the other participants.

\begin{listing}
\begin{minted}{C}
#define SUB_SLOT_POW2  3
#define SUB_SLOT_CNT   (1<<SUB_SLOT_POW2)   // 8 sub-slots per slot
#define SUB_SLOT_MASK  (SUB_SLOT_CNT-1)     // 0x07
#define SLOT_POW2      4
#define SLOT_CNT       (1<<SLOT_POW2)        // 16 slots per cycle
#define SLOT_MASK      (SLOT_CNT-1)          // 0x0F
#define SLOT_SHIFT     (SUB_SLOT_POW2)       // 3
\end{minted}
\caption{Slot constants from \texttt{system\_definitions.h}}
\label{cycle_defines}
\end{listing}

\subsubsection*{State}
The entire state is held in the structure \verb|cycle_t| (Listing
\ref{cycle_t}). The central field is the monotonically increasing counter
\verb|subSlot| (range $0 \dots 127$), in which slot and sub-slot index are
\emph{bit-packed}: the lower \verb|SUB_SLOT_POW2| bits hold the sub-slot, the
\verb|SLOT_POW2| bits above hold the slot. The two macros from Listing
\ref{cycle_macros} extract the two parts so that no computation is needed
elsewhere:
\begin{align*}
\mathtt{ACT\_SUB\_SLOT} &= \mathtt{subSlot}\ \&\ \mathtt{SUB\_SLOT\_MASK} \\
\mathtt{ACT\_SLOT}      &= (\mathtt{subSlot} \gg \mathtt{SLOT\_SHIFT})\ \&\ \mathtt{SLOT\_MASK}
\end{align*}
The fields \verb|actSlot| and \verb|sSlot| are cached copies of the current
slot and sub-slot, updated on every advance so that they are available without
a renewed bit operation. \verb|cycle| counts the fully completed cycles
($16$ bit, wraps at $65535$) and \verb|init| marks successful initialisation.

\begin{listing}
\begin{minted}{C}
typedef struct cycle_s {
    volatile int8_t subSlot; // current sub-slot (slot+sub-slot packed)
    int8_t          actSlot; // cached slot
    int8_t          sSlot;   // cached sub-slot
    uint16_t        cycle;   // number of completed cycles
    bool            init;    // true after cycle_init()
} cycle_t;
\end{minted}
\caption{Definition of the \texttt{cycle\_t} structure}
\label{cycle_t}
\end{listing}

\begin{listing}
\begin{minted}{C}
#define ACT_SUB_SLOT(_cycle)  (((_cycle)->subSlot) & SUB_SLOT_MASK)
#define ACT_SLOT(_cycle)      (((_cycle)->subSlot >> SLOT_SHIFT) & SLOT_MASK)
\end{minted}
\caption{Macros for extracting slot and sub-slot}
\label{cycle_macros}
\end{listing}

\subsubsection*{API}
\begin{description}
\item[\texttt{cycle\_init(cycle)}] Sets \verb|init| and then calls
  \verb|cycle_reset()|. Returns: \verb|em_msg| (\verb|EM_OK|, or \verb|EM_ERR|
  on a null pointer).
\item[\texttt{cycle\_reset(cycle)}] Resets \verb|subSlot| and \verb|cycle| to
  $0$. Requires a previously initialised structure. Returns: \verb|em_msg|
  (\verb|EM_ERR| if not initialised).
\item[\texttt{cycle\_set\_slot(cycle, slot)}] Positions the counter at the
  start of the given slot ($\mathtt{subSlot} = \mathtt{slot} \cdot
  \mathtt{SUB\_SLOT\_CNT}$). Only valid slots are accepted; the slot is checked
  beforehand with \verb|cycle_check_slot()|. Returns: \verb|em_msg|
  (\verb|EM_OK| for a valid slot, \verb|EM_ERR| on a null pointer or invalid
  slot).
\item[\texttt{cycle\_check\_slot(slot)}] Validates a slot: odd values in the
  range $0 < \mathtt{slot} \le \mathtt{SLOT\_CNT}$ are allowed. Returns:
  \verb|int8_t| -- the slot or $-1$.
\item[\texttt{cycle\_check(cycle, rxSlot, ss)}] Checks whether a frame received
  from slot \verb|rxSlot| fits the current cycle position. The result is true
  if the current slot equals \verb|rxSlot|, or if \verb|rxSlot| is the
  immediately \emph{following} slot and the current sub-slot is smaller than
  the tolerance \verb|ss|, or if \verb|rxSlot| is the immediately
  \emph{preceding} slot and fewer than \verb|ss| sub-slots remain until the end
  of the slot (slot indices taken modulo \verb|SLOT_CNT|). The tolerance window
  \verb|ss| therefore also admits frames from the directly neighbouring slots.
  Returns: \verb|bool| (\verb|false| on a null pointer or invalid
  \verb|rxSlot|).
\item[\texttt{cycle\_difference(cycle, rxSlot)}] Returns the signed distance in
  sub-slots from the current position \verb|subSlot| to the \emph{window} of
  slot \verb|rxSlot|. That window spans \verb|SUB_SLOT_CNT| sub-slots with the
  lower corner $\mathtt{rxSlot} \cdot \mathtt{SUB\_SLOT\_CNT}$ and the upper
  corner $\mathtt{rxSlot} \cdot \mathtt{SUB\_SLOT\_CNT} + \mathtt{SUB\_SLOT\_CNT} - 1$.
  Inside the window the distance is $0$; below the lower corner it is the
  (negative) distance to that corner, above the upper corner the (positive)
  distance past it. It is measured the shortest way around the ring of
  $\mathtt{SUB\_SLOT\_CNT} \cdot \mathtt{SLOT\_CNT} = 128$ sub-slots, so
  \verb|rxSlot|~$0$ and \verb|rxSlot|~$=\mathtt{SLOT\_CNT}$ yield the same
  result. With the signed distance to the lower corner
  \[
    d = \big((\mathtt{subSlot} - \mathtt{rxSlot} \cdot
    \mathtt{SUB\_SLOT\_CNT} + 64) \bmod 128\big) - 64
  \]
  the result is
  \[
    \mathtt{diff} =
    \begin{cases}
      d & d < 0 \quad(\text{below the lower corner}),\\
      0 & 0 \le d < \mathtt{SUB\_SLOT\_CNT} \quad(\text{inside the window}),\\
      d - (\mathtt{SUB\_SLOT\_CNT}-1) & d \ge \mathtt{SUB\_SLOT\_CNT}
        \quad(\text{above the upper corner}).
    \end{cases}
  \]
  Returns: \verb|int8_t| ($0$ on a null pointer or uninitialised structure).
\item[\texttt{cycle\_string(cycle)}] Returns a compact representation in the
  format \verb|(c:<cycle>, <slot>, <subSlot>)| in a static buffer. Returns:
  \verb|char *| (\verb|NULL| on a null pointer).
\item[\texttt{cycle\_print(cycle, title)}] Prints all fields line by line
  (debug aid, cf.\ the \verb|xx_print| convention in the introduction).
  Returns: \verb|em_msg| (\verb|EM_ERR| on a null pointer or uninitialised
  structure, otherwise \verb|EM_OK|).
\end{description}

\subsubsection*{Advancing the cycle}
{\sloppy\emergencystretch=3em
The counter is advanced solely through \verb|cycle_increment()|, typically
once per sub-slot tick. The function is coupled to the synchronisation state
machine (\verb|system_state_e|):

\begin{itemize}
\item In state \verb|SYNCHRONIZE| the transition to \verb|SYNCHRONIZE_READY| is
  triggered and the internal flag \verb|is_set| is set; from this point on the
  counter runs.
\item As long as \verb|is_set| holds, \verb|subSlot| is incremented and folded
  back modulo $\mathtt{SUB\_SLOT\_CNT} \cdot \mathtt{SLOT\_CNT}$ ($=128$);
  \verb|actSlot| and \verb|sSlot| are derived from the new value.
\item In states \verb|SYNCHRONIZE_DOING|, \verb|SYNCHRONIZE_READY| and
  \verb|SYNCHRONIZE_ERROR| a slot change is detected. On the transition to
  slot $0$ the cycle counter \verb|cycle| is incremented exactly once per pass
  (debounced via the flag \verb|cycle_once|).
\end{itemize}
\par}

\noindent Since \verb|cycle_increment()| advances first and only then reports
the fields, the first call after a reset already yields the next sub-slot. If
\verb|OPTION_SHOW_TIMING| is active, status LEDs are additionally toggled for
run-time measurement.
